Este trabalho partiu da pergunta formulada na Seção~\ref{sec:problema}: de que maneira a
adoção do processamento em borda ou em nuvem impacta a eficiência energética, a acurácia, a
complexidade de desenvolvimento de firmware e o desempenho computacional de um sistema
\gls{iot} voltado à coleta de dados de medidores analógicos. Para respondê-la, os dois
fluxos foram implementados sobre o mesmo nó sensor, baseado no \gls{esp32cam} e no módulo
Radioenge RD49C, e sobre o mesmo enlace \gls{lorawan}, e comparados em bancada.

A resposta obtida é que, em um enlace \gls{lorawan} de baixa taxa, o volume de dados
transmitido pesa mais do que a capacidade de processamento. Embora o \gls{esp32cam} execute a
inferência cerca de 2.800 vezes mais devagar que o servidor, o fluxo em borda, que transmite
cerca de 450 vezes menos dados por leitura, teve latência de 10 a 54 vezes menor, conforme
executasse a localização da janela ou usasse a ROI calibrada, e consumo por leitura de 5,6
a 16,7 vezes menor, conforme a configuração. A acurácia não diferenciou os fluxos, que executam
os mesmos modelos. O custo do processamento local apareceu no tempo de processamento e na
memória exigidos do nó e na maior complexidade de desenvolvimento e de manutenção do
firmware.

O objetivo geral foi atingido: a viabilidade técnica, a acurácia e o consumo energético dos
dois fluxos foram comparados com medições no protótipo. Quanto aos objetivos específicos:

\begin{enumerate}
    \item os dois fluxos foram implementados no mesmo hardware e operaram de forma autônoma,
    em ciclos de captura, envio e repouso em \textit{deep sleep};
    \item o enlace foi configurado no plano AU915, em DR6 (SF8 com 500~kHz) e canal fixo,
    enquadrado no item 10.3 do Ato n\textordmasculine~14.448 da ANATEL. Para reduzir o
    consumo, o nó dorme entre as leituras e desliga a câmera logo após a captura, o que
    reduziu em 30\% a carga por leitura do fluxo em nuvem. Com o módulo de rádio configurado
    em classe A, a corrente em repouso caiu de cerca de 24 para cerca de 6~mA. Esse objetivo
    foi atingido em parte, pois esses cerca de 6~mA, ainda não separados entre o módulo e a
    placa, dominam a autonomia em intervalos longos;
    \item os dois fluxos foram avaliados segundo as métricas previstas. Na avaliação em
    computador, com a janela de dígitos marcada à mão, os dois tiveram a mesma precisão no
    conjunto de teste comum: 46\% das leituras e 85\% dos dígitos corretos; com o quadro
    endireitado, recurso implementado apenas no servidor, a leitura chegou a 82\% e 92\%. No
    hidrômetro físico, com a câmera fixa, a ROI calibrada e o recorte da face da roda, os dois
    fluxos leram 25 das 30 fotos (83\%) e 97\% dos dígitos. A latência de ponta a ponta foi
    de 189,1~s no fluxo em nuvem e, no fluxo em borda, de 18,84~s com a localização da janela
    e de 3,50~s com a ROI calibrada; o consumo foi de 4,64~mAh por leitura no fluxo em nuvem
    e de 0,824~mAh no fluxo em borda com a localização e o rádio em classe C, e de 2,91 e
    0,175~mAh com a ROI e o rádio em classe A; e o volume
    transmitido foi de 9~bytes por leitura no fluxo em borda, contra cerca de 4~KB de imagem
    no fluxo em nuvem.
\end{enumerate}

Quanto às hipóteses da Seção~\ref{subsec:hipotese}, cujos resultados estão resumidos na
Tabela~\ref{tab:sintese_hipoteses}:

\begin{itemize}
    \item a H1, avaliada pela comparação entre as versões em ponto flutuante e em INT8 do
    mesmo classificador, foi confirmada no sentido, mas a redução de acurácia foi de apenas
    0,08 ponto percentual (95,38\% contra 95,46\%) e não teve significância estatística
    ($p = 0{,}29$); no hidrômetro físico, as duas versões leram as mesmas 25 das 30 fotos, de
    modo que a quantização não penalizou o reconhecimento;
    \item a H2 foi confirmada: no nó, uma leitura do fluxo em borda exigiu 17.568~ms de
    processamento e 1,65~MB de PSRAM, contra cerca de 11~ms de compressão JPEG no fluxo em
    nuvem;
    \item a H3 foi confirmada: nas duas campanhas, o fluxo em nuvem consumiu mais por leitura
    e teve autonomia inferior. Na configuração adotada, com a ROI e o rádio em classe A, uma
    leitura custou 2,91~mAh no fluxo em nuvem, contra 0,175~mAh no fluxo em borda. A magnitude
    da diferença na autonomia, porém, depende da corrente em repouso e do intervalo entre
    leituras: para uma bateria hipotética de 2.000~mAh, o fluxo em borda teria 12,1 dias de
    autonomia, contra 3,8 dias do fluxo em nuvem, com leituras a cada 10~minutos, e 13,2
    contra 9,5 dias com leituras a cada hora; com o rádio em classe C, as autonomias com
    leituras horárias eram de apenas 3,4 e 3,1 dias.
\end{itemize}

A entrega dos dados, avaliada de forma complementar, não diferenciou os fluxos: por leitura,
o fluxo em borda entregou 79,4\% e o fluxo em nuvem, 66,7\%, sem significância estatística
($p = 0{,}27$).
